\documentclass[12pt]{article}
\usepackage[utf8]{inputenc}
\usepackage[T1]{fontenc}
\usepackage{hyperref}
\title{Tombulu Language Resources for K-12 Curriculum Development}
\author{Compiled by Claude Sonnet 5 (claude-sonnet-5) for Andr\'as Kornai}
\date{}
\begin{document}
\maketitle

\section{Introduction}
Tombulu is a Minahasan Austronesian language of North Sulawesi, spoken around Tomohon. The reference description gives about 6,000 speakers, while one source cites about 19,000 from the 2010 census, so the figures differ. Little is available online. This report is based on one search round.

\section{Primary Educational Resources}
No school materials or school programmes were found.

\section{Linguistic and Reference Materials}
\begin{itemize}
\item Overview and phonology: \url{https://en.wikipedia.org/wiki/Tombulu_language}
\item Writing and language overview: \url{https://omniglot.com/writing/tombulu.htm}
\item Academic article: \url{https://journal.iain-manado.ac.id/index.php/JII/article/download/561/465}
\item Academic article: \url{https://ejurnal.unima.ac.id/index.php/e-clue/article/download/2971/1272/9625}
\end{itemize}
Print works cited by a Wikipedia mirror are Sahulata (1993), \emph{Struktur bahasa Tombulu}, and Makalew-Palar et al. (1994), \emph{Fonologi Bahasa Tombulu}. I did not locate copies. I could not confirm the publisher or title of an Indonesian--Tombulu dictionary said to be released in 2013.

\section{Religious and Community Resources}
A Tombulu New Testament was reportedly released in 2018, per Wikipedia-derived text. I did not find a URL to the edition itself.

\section{Digital and Audio Resources}
None found.

\section{Curriculum Development Context}
Very little is available. A curriculum would first require documentation work with the Tomohon community. The Indonesian-language academic sources above are the best starting points.
\end{document}